\chapter{HILBERT SPACES}

For the most elementary facts about Hilbert spaces I can think of no better introduction than the first two chapters
of Halmos's classic text~\cite{Halmos:1951}.  Another good introduction is~\cite{Young:1988}.  You may also find
chapters 3 and 9 of~\cite{Kreyszig:1978} and chapter 4 of~\cite{Ha:2006} helpful.  For a start towards more serious
explorations into the mysteries of Hilbert space look at~\cite{Conway:1990}, \cite{Conway:2000}, and~\cite{Helmberg:1969}.


\section{Definition and Examples}

A Hilbert space is a \emph{complete inner product space}.  In more detail:
\begin{defn} In proposition~\ref{00015025} we showed how an inner product on a vector space induces a norm on the space and in
proposition~\ref{0001503} how a norm in turn induces a metric.  If an inner product space is complete with respect to this metric
it is a
  \index{Hilbert!space}%
  \index{space!Hilbert}%
\df{Hilbert space}.  Similarly, a
  \index{Banach!space}%
  \index{space!Banach}%
\df{Banach space} is a complete normed linear space and a
  \index{Banach!algebra}%
  \index{algebra!Banach}%
\df{Banach algebra} is a complete normed algebra.
\end{defn}

\begin{exam} Under the inner product defined in
 \index{K@$\K^n$!as a Hilbert space}%
 \index{Hilbert!space!$\K^n$ as a}%
example~\ref{000150012} $\K^n$ is a Hilbert space.
\end{exam}

\begin{exam}\label{exam150013} Under the inner products in
 \index{l@$l_2 = l_2(\N)$!as a Hilbert space}%
 \index{l@$l_2(\Z)$!as a Hilbert space}%
 \index{Hilbert!space!$l_2 = l_2(\N)$, $l_2(\Z)$ as a}%
example~\ref{000150013} $l_2 = l_2(\N)$ and $l_2(\Z)$ are Hilbert spaces.
\end{exam}

\begin{exam} Under the inner product defined in example~\ref{000150014} $\fml C([a,b])$ is \emph{not} a Hilbert space.
\end{exam}

\begin{exam}\label{exam_CX} Let $\fml C(X) = \fml C(X,\C)$ be the family of all continuous complex valued functions on a compact
Hausdorff space~$X$.  Under the
 \index{uniform!norm!on $\fml C(X)$}%
 \index{norm!uniform}%
 \index{C@$\fml C(X)$!as a Banach space}%
 \index{C@$\fml C(X)$!continuous complex valued functions on~$X$}%
\df{uniform norm}
   \[ \norm f_u := \sup\{\abs{f(x)}\colon 0 \le x \le 1\} \]
$\fml C(X)$ is a Banach space.  If the space $X$ is not compact we can make the set $\fml C_b(X)$
 \index{C@$\fml C_b(X)$!as a Banach space}%
of all \emph{bounded} continuous complex valued functions on $X$ into a Banach space in the same way.
When the scalar field is $\R$ we will
 \index{C@$\fml C(X,\R)$!continuous real valued functions on~$X$}%
 \index{C@$\fml C(X,\R)$!as a Banach space}%
write $\fml C(X,\R)$ for the family of continuous real valued functions on~$X$.  It too is a Banach space.
\end{exam}

\begin{exam}\label{000213} The inner product space $l_c$ of all sequences $(a_n)$ of complex numbers which are eventually zero
(see example~\ref{exam13122}) is not a Hilbert space.
\end{exam}

\begin{exam}\label{00022} Let $\mu$ be a positive measure on a $\sigma$-algebra $\sfml A$ of
subsets of a set~$S$.  If you are unacquainted with general measure theory, a casual introduction to Lebesgue measure
on the real line should suffice for examples which appear in these notes.  A complex valued function $f$ on $S$ is
 \index{measurable}%
\df{measurable} if the inverse image under $f$ of every Borel set (equivalently, every open
set) in $\C$ belongs to~$\sfml A$. We define an equivalence relation $\sim$ on the family of
measurable complex valued functions by setting $f \sim g$ whenever $f$ and $g$ differ on a set
of measure zero, that is, whenever $\mu(\{x \in S\colon f(x) \ne g(x)\}) = 0$.  We adopt
conventional notation and denote the equivalence class containing $f$ by $f$ itself (rather
than something more logical such as~$[f]$).  We denote the family of (equivalence classes of)
measurable complex valued functions on $S$ by $\fml M(S,\mu)$. A function $f \in \fml M(S,\mu)$ is
 \index{square integrable}%
 \index{M@$\fml M(S,\mu)$ (equivalence classes of measurable functions)}%
 \index{L@$\fml L_2(S,\mu)$!square integrable functions}%
\df{square integrable} if $\int_S \abs{f(x)}^2\,d\mu(x) < \infty$.  We denote the family of
(equivalence classes of) square integrable functions on $S$ by~$\lfs 2(S,\mu)$.
 \index{L@$\lfs 2(S,\mu)$!classes of square integrable functions}%
 \index{L@$\lfs 2(S,\mu)$!as a Hilbert space}%
 \index{Hilbert!space!$\lfs 2(S,\mu)$ as a}%
For every $f$, $g \in \lfs 2(S,\mu)$ define
   \[ \langle f,g \rangle := \int_S f(x) \conj{g(x)}\,d\mu(x)\,. \]
With this inner product (and the obvious pointwise vector space operations) $\lfs 2(S,\mu)$ is a Hilbert space.
\end{exam}

\begin{exam}\label{exam_L1S} Notation as in the preceding example.  A function $f \in \fml M(S,\mu)$ is
 \index{integrable}%
 \index{L@$\fml L_1(S,\mu)$!integrable functions}%
\df{integrable} if $\int_S \abs{f(x)}\,d\mu(x) < \infty$.  We denote the family of
(equivalence classes of) integrable functions on $S$ by~$\lfs 1(S,\mu)$.
 \index{L@$\lfs 1(S,\mu)$!classes of integrable functions}%
 \index{L@$\lfs 1(S,\mu)$!as a Banach space}%
 \index{Banach!space!$\lfs 1(S,\mu)$ as a}%
For every $f \in \lfs 1(S,\mu)$ define
    \[ {\norm f}_1 := \int_S \abs{f(x)}\,d\mu(x)\,. \]
With this norm (and the obvious pointwise vector space operations) $\lfs 1(S,\mu)$ is a Banach space.
\end{exam}

\begin{defn} If $J$ is any interval in $\R$ a real (or complex) valued function $f$ on $J$ is
 \index{absolute!continuity!of a function}%
 \index{continuity!absolute!of a function}%
 \index{absolutely!continuous}%
\df{absolutely continuous} if it satisfies the following condition:
  \quote for every $\epsilon > 0$ there exists $\delta > 0$ such
 that if $(a_1,b_1), \dots,(a_n,b_n)$ \\
 are disjoint nonempty subintervals of $J$ with $\sum_{k=1}^n(b_k - a_k) < \delta$, \\
 then $\sum_{k=1}^n\abs{f(b_k) - f(a_k)}~<~\epsilon$.
  \endquote
\end{defn}

\begin{exam}\label{exam_Hsp_abs_cont} Let $\lambda$ be Lebesgue measure on the interval $[0,1]$ and $H$ be the set of all
absolutely continuous functions on $[0,1]$ such that $f\,'$ belongs to $\lfs 2([0,1],\lambda)$ and $f(0) = 0$.
 \index{Hilbert!space!of absolutely continuous functions}%
 \index{absolutely!continuous functions!Hilbert space of}%
For $f$ and $g$ in $H$ define
   \[ \langle f,g \rangle = \int_0^1 f\,'(t)\clo{g\,'(t)}\,dt. \]
This is an inner product on $H$ under which $H$ becomes a Hilbert space.
\end{exam}

\begin{exam} If $H$ and $K$ are Hilbert spaces, then their (external orthogonal) direct sum $H \oplus K$
 \index{Hilbert!space!direct sum of Hilbert spaces as a}%
 \index{direct!sum!of Hilbert spaces}%
(as defined in~\ref{0001504}) is also a Hilbert space.
\end{exam}


















\section{Unordered Summation}

Defining \emph{unordered summation} is no more difficult in normed linear spaces than in~$\R$.  We generalize
definition~\ref{def_summable_R}.

\begin{defn} Let $V$ be a normed linear space and $A \subseteq V$.
For every $F \in \fin A$ define
  \[ \sbsb sF = \sum F\]
where $\sum F$ is the sum of the (finitely many) vectors belonging to~$F$.
Then $s = \bigl(\sbsb sF\bigr)_{F \in \fin A}$ is a net in~$V$. If this net converges, the set
$A$ is said to be
 \index{summable!set of vectors}%
\df{summable}; the limit of the net is the
 \index{sum!of a summable set}%
 \index{<sum@$\sum A$ (sum of a set of vectors)}%
\df{sum} of $A$ and is denoted by~$\sum A$. If $\{\norm a\colon a \in A\}$ is summable, we say that the set $A$ is
 \index{absolutely!summable}%
 \index{summable!absolutely}%
\df{absolutely summable}.

To accommodate sums in which the same vector appears more than once, we need to make use of indexed families of vectors.  They
require a slightly different approach. Suppose, for example, that $A = \bigl(x_i\bigr)_{i \in I}$ where $I$ is
an arbitrary index set. Then for each $F \in \fin I$
  \[ \sbsb sF = \sum_{i \in F}x_i . \]
As above $s$ is a net in~$V$. If it converges $\bigl(x_i\bigr)_{i \in I}$ is
 \index{summable!indexed family of vectors}%
summable and its limit is the
 \index{sum!of an indexed family}%
 \index{<sum@$\sum_{i \in I} x_i$ (sum of an indexed family)}%
sum of the indexed family, and is denoted by $\sum_{i \in I} x_i$.  An alternative way of saying that $\bigl(x_i\bigr)_{i \in I}$
is summable is to say that the ``series'' $\sum_{i \in I}x_i$
 \index{convergence!of infinite series}%
 \index{sum!of an infinite series}%
\df{converges} or that the ``sum'' $\sum_{i \in I}x_i$ \df{exists}.  An indexed family $\bigl(x_i\bigr)_{i \in I}$ is
 \index{absolutely!summable}%
 \index{summable!absolutely}%
\df{absolutely summable} if the indexed family $\bigl(\norm{x_i}\bigr)_{i \in I}$ is summable.  An alternative way of saying
that $\bigl(x_i\bigr)_{i \in I}$ is absolutely summable is to say that the ``series'' $\sum_{i \in I}x_i$ is
 \index{convergence!absolute}%
 \index{absolute!convergence}%
\df{absolutely convergent}.

Sequences, as usual, are treated as a separate case.  Let $(a_k)$ be a sequence in~$V$.  If the infinite series
$\sum_{k=1}^\infty a_k$ (that is, the sequence of
 \index{partial!sum}%
 \index{sum!partial}%
 \index{infinite series!partial sum of an}%
 \index{series!infinite}%
partial sums $s_n = \sum_{k=1}^n a_k$) converges to a
vector $b$ in $V$, then we say that the sequence $(a_k)$ is
 \index{infinite series!in a normed linear space}%
 \index{summable!sequence}%
 \index{sequence!summable}%
\df{summable} or, equivalently, that the series $\sum_{k=1}^\infty a_k$ is a
 \index{convergence!of a series}%
 \index{series!convergent}%
\df{convergent series}.  The vector $b$ is called the
 \index{sum!of an infinite series}%
 \index{series!sum of a}%
 \index{infinite series!sum of an}%
\df{sum} of the series $\sum_{k=1}^\infty a_k$ and we write
  \[ \sum_{k=1}^\infty a_k = b\,. \]
It is clear that a necessary and sufficient condition for a series $\sum_{k=1}^\infty a_k$ to
be convergent or, equivalently, for the sequence $(a_k)$ to be summable, is that there exist a
vector $b$ in $V$ such that
  \begin{equation}\label{cond_ser_conv}
     \biggnorm{b - \sum_{k=1}^n a_k} \sto 0 \qquad \text{as} \qquad n \sto \infty.
  \end{equation}
\end{defn}

\begin{prop}\label{prop_abs_sum_implies_sum} A normed linear space $V$ is complete if and only if every absolutely
summable sequence in $V$ is summable.
\end{prop}

\begin{proof}[\emph{Hint for proof}] Let $(x_n)$ be a Cauchy sequence in a normed linear space in which every absolutely
summable sequence is summable.  Find a subsequence $\bigl(x_{n_k}\bigr)$ of $(x_n)$ such that $\norm{x_n - x_m} < 2^{-k}$
whenever $m$, $n \ge n_k$.  If we let $y_1 = x_{n_1}$ and $y_j = x_{n_j} - x_{n_{j-i}}$ for $j > 1$, then $(y_j)$ is
absolutely summable.   \ns
\end{proof}

\begin{prop} If $\bigl(x_i\bigr)_{i \in I}$ and $\bigl(y_i\bigr)_{i \in I}$ are
summable indexed families in a normed linear space and $\alpha$ is a scalar, then
$\bigl(\alpha x_i\bigr)_{i \in I}$ and $\bigl(x_i + y_i\bigr)_{i \in I}$
are summable; and
  \[ \sum_{i \in I} \alpha x_i = \alpha \, \sum_{i \in I} x_i \]
and
  \[ \sum_{i \in I} x_i + \sum_{i \in I} y_i =
                         \sum_{i \in I} (x_i + y_i). \]
\end{prop}

\begin{rem} One frequently finds some version of the
following ``proposition'' in textbooks.
\vskip 10 pt
\begin{quote}
If $\sum_{i \in I}x_i = u$ and $\sum_{i \in I}x_i = v$ in a
Hilbert (or Banach) space, then $u = v$.
\end{quote}

 \vskip 10 pt
\noindent (I once saw a 6 line proof of this.) Of which result is
this a trivial consequence?
\end{rem}

\begin{prop} If $\bigl(x_i\bigr)_{i \in I}$ is a summable indexed family of vectors in a Hilbert space $H$ and
$y$ is a vector in $H$, then $\bigl(\langle x_i,y \rangle\bigr)_{i \in I}$ is a summable indexed family of scalars and
   \[ \sum_{i \in I}\langle x_i,y \rangle  = \biggl < \sum_{i \in I} x_i, y \biggr >. \]
\end{prop}

The following is a generalization of example~\ref{exam150013}.

\begin{exam} Let $I$ be an arbitrary nonempty set and $l^2(I)$ be the set of all complex valued functions $x$ on $I$ such that
$\sum_{i \in I}\abs{x_i}^2 < \infty$. For all $x$, $y \in l^2(I)$ define
    \[ \langle x,y \rangle = \sum_{i \in I}x(i)\conj{y(i)}.\]
  \begin{enumerate}
    \item If $x \in l^2(I)$, then $x(i) = 0$ for all but countably many $i$.
    \item The map $(x,y) \mapsto \langle x,y \rangle$ is an inner product on the vector space $l^2(I)$.
    \item The space $l^2(I)$ is a Hilbert space.
  \end{enumerate}
\end{exam}

\begin{prop}\label{prop_sum_conv} Let $(x_n)$ be a sequence in a Hilbert space.  If the sequence regarded as an indexed family
is summable, then the infinite series $\sum_{k=1}^\infty x_k$ converges.
\end{prop}

\begin{exam} The converse of proposition~\ref{prop_sum_conv} does not hold.
\end{exam}

The direct sum of two Hilbert spaces is a Hilbert space.

\begin{prop} If $H$ and $K$ are Hilbert spaces, then their (orthogonal) direct sum (as defined in~\ref{0001504})
is a Hilbert space.
\end{prop}

We can also take the product of more than two Hilbert spaces---even an infinite family of them.

\begin{prop} Let $\bigl(H_i\bigr)_{i \in I}$ be an indexed family of Hilbert spaces.  A function $x\colon I \sto \bigcup H_i$
belongs to $H = \D\bigoplus_{i \in I}H_i$ if $x(i) \in H_i$ for every $i$ and if $\sum_{i \in I}\norm{x(i)}^2 < \infty$.
Defining $\langle \hphantom x, \hphantom y \rangle$ on $H$ by $\langle x, y \rangle := \sum_{i \in I} \langle
x(i), y(i) \rangle$ makes $H$ into an Hilbert space.
\end{prop}

\begin{defn} The Hilbert space defined in the preceding proposition is the
 \index{external!othogonal direct sum}%
 \index{orthogonal!direct sum!of Hilbert spaces}%
 \index{direct!sum!of Hilbert spaces}%
 \index{Hilbert!space!direct sum}%
 \index{sum!direct}%
\df{(external orthogonal) direct sum} of the Hilbert spaces~$H_i$.
\end{defn}

\begin{exer}  Is the direct sum defined above (with appropriate morphisms) a product in the category $\cat{HSp}$
of Hilbert spaces and bounded linear maps?  Is it a coproduct?
\end{exer}


























\section{Hilbert space Geometry}

\begin{conv}  In the context of Hilbert spaces the word ``subspace'' will always mean
 \index{conventions!subspaces of Hilbert spaces are closed}%
 \index{subspace!of a Hilbert space}%
\emph{closed vector subspace}.  The reason for this is that we want a subspace of a Hilbert space to be a Hilbert space in its own right; that is to say,
a subspace of a Hilbert space should be a subobject in the category of Hilbert spaces (and appropriate morphisms).  To indicate that $M$ is a subspace of $H$ we write
 \index{<binrelle@$\preccurlyeq$ (subspace of a Banach or Hilbert space)}%
$M \preccurlyeq H$.  A (not necessarily closed) vector subspace of a Hilbert space is also
 \index{linear!subspace}%
 \index{vector!subspace}%
 \index{linear!manifold}%
 \index{manifold!linear}%
called a \emph{linear subspace} or a \emph{linear manifold}.
\end{conv}

\begin{defn} Let $A$ be a nonempty subset of a Banach (or Hilbert) space~$B$.  We define the
 \index{closed!linear span}%
 \index{linear!span!closed}%
 \index{<unops@$\bigvee A$ (closed linear span of~$A$)}%
\df{closed linear span} of $A$ (denoted by $\bigvee A$) to be the smallest subspace of $B$ containing~$A$; that is, the
intersection of all subspaces of $B$ which contain~$A$.
\end{defn}

As with definitions \ref{smplx005def}, \ref{def_interior}, and \ref{def_closure}, we need to show that the preceding one makes sense.

\begin{prop} The preceding definition makes sense.  Furthermore, a vector is in the closed linear span of $A$ if and only if it belongs
to the closure of the linear span of~$A$.
\end{prop}

\begin{thm}[Minimizing Vector Theorem] If $C$ is a nonempty closed convex subset of a
 \index{minimizing vector theorem}%
Hilbert space $H$ and $a \in C^c$, then there exists a unique $b \in C$ such that
$\norm{b - a} \le \norm{x - a}$ for every $x \in C$.
 \[\xymatrix{
     &&&&&{}  \\
     &&&& x   \\
     a\ar[rrrru]\ar[rrr] &&& b  \\
     &&&&&C   \\
     &&&&{}\ar@{-}@/^3.5pc/[uuuu]       }\]
\end{thm}

\begin{proof}[\emph{Hint for proof}]  For the existence part let $d = \inf\{\norm{x - a}\colon x \in C\}$.  Choose a sequence
$(y_n)$ of vectors in $C$ such that $\norm{y_n - a} \sto d$.  To prove that the sequence $(y_n)$ is Cauchy, apply the
\emph{parallelogram law}~\ref{parallelogram_law} to the vectors $y_m - a$ and $y_n - a$ for $m$, $n \in \N$.  The uniqueness
part is similar: apply the \emph{parallelogram law} to the vectors $b - a$ and $c - a$ where $b$ and $c$ are two vectors
satisfying the conclusion of the theorem.  \ns
\end{proof}

\begin{exam} The vector space $\R^2$ under the uniform metric is a Banach space. To see
that in this space the \emph{minimizing vector theorem} does not hold take $C$ to be the
closed unit ball about the origin and $a$ to be the point~$(2,0)$.
\end{exam}

\begin{exam} The sets
   \[ C_1 = \biggl\{f \in \fml C([0,1])\colon \int_0^{1/2}f - \int_{1/2}^1f = 1\biggr\} \qquad \text{and} \qquad
                              C_2 = \biggl\{f \in \lfs 1([0,1],\lambda)\colon \int_0^1f = 1\biggr\} \]
are examples that show that neither the existence nor the uniqueness claims of the \emph{minimizing vector theorem}
necessarily holds in a Banach space. (The symbol $\lambda$ denotes Lebesgue measure.)
\end{exam}

\begin{thm}[Vector Decomposition Theorem]\label{thm_vec_decomp} Let $H$ be a Hilbert space and
 \index{vector!decomposition theorem}%
$M$ be a (closed linear) subspace of~$H$.
Then for every $x \in H$ there exist unique vectors $y \in M$ and $z \in M^\perp$ such that $x = y + z$.
\end{thm}

\begin{proof}[\emph{Hint for proof}] Assume $x \notin M$ (otherwise the result is trivial). Explain how we
know that there exists a unique vector $y \in M$ such that
   \begin{equation} \norm{x - y} \le \norm{x - m} \end{equation}
for all $m \in M$.  Explain why it is enough, for the existence part of the proof, to prove
that $x - y$ is perpendicular to every \emph{unit} vector $m \in M$.  Explain why it is true
that for every unit vector $m \in M$
   \begin{equation} \norm{x - y}^2 \le \norm{x - (y + \lambda m)}^2 \end{equation}
where $\lambda := \langle x - y, m \rangle$.  \ns
\end{proof}

\begin{cor} If $M$ is a subspace of a Hilbert space~$H$, then $H = M \oplus M^\perp$.
\end{cor}

\begin{exam}The preceding result says that every subspace of a Hilbert space is a direct
summand. This result need not be true if $M$ is assumed to be just a linear subspace of the space.
For example, notice that $M = l_c$ (see example~\ref{000213}) is a linear subspace of the
Hilbert space $l_2$ (see example~\ref{exam150013}) but $M^\perp = \{\vc 0\}$.
\end{exam}

\begin{cor} A vector subspace $A$ of a Hilbert space $H$ is dense in $H$ if and only if $A^\perp = \{0\}$.
\end{cor}

\begin{prop} Every proper subspace of a Hilbert space has empty interior.
\end{prop}

\begin{prop} If a Hamel basis for a Hilbert space is not finite, then it is uncountable.
\end{prop}

\begin{proof}[\emph{Hint for proof}] Suppose the space has a countably infinite basis $\{e^1, e^2, e^3, \dots \}$.
For each $n \in \N$ let $M_n = \spn\{e^1, e^2, \dots, e^n\}$.  Apply the \emph{Baire category theorem} (see~\cite{Erdman:2007},
theorem 29.3.20) to $\bigcup_{n=1}^\infty M_n$ keeping in mind the preceding proposition.   \ns
\end{proof}

\begin{prop}\label{00023} Let $H$ be a Hilbert space. Then the following hold:
 \begin{enumerate}
  \item if $A \subseteq H$, then $A \subseteq A^{\perp\perp}$;
  \item if $A \subseteq B \subseteq H$, then $B^\perp \subseteq A^\perp$;
  \item if $M$ is a subspace of $H$, then $M = M^{\perp\perp}$; and
  \item if $A \subseteq H$, then $\bigvee A = A^{\perp\perp}$.
 \end{enumerate}
\end{prop}

\begin{prop} Let $M$ and $N$ be subspaces of a Hilbert space.  Then
 \begin{enumerate}
  \item  $(M + N)^\perp = (M \cup N)^\perp =  M^\perp \cap N^\perp$, and
  \item  $(M \cap N)^\perp = M^\perp + N^\perp$.
 \end{enumerate}
\end{prop}

\begin{prop} Let $V$ be a normed linear space.  There exists an inner product on $V$ which induces the norm on $V$
if and only if the norm satisfies the \emph{parallelogram law}~\ref{parallelogram_law}.
\end{prop}

\begin{proof}[\emph{Hint for proof}]  We say that an inner product $\langle \,\cdot\,\,,\,\cdot\, \rangle$
on $V$ induces the norm $\norm{\,\cdot\,}$ if $\norm x = \sqrt{\langle x,x \rangle}$ holds for every $x$ in $V$.)
To prove that if a norm satisfies the \emph{parallelogram law} then it is induced by an inner product,
use the equation in the \emph{polarization identity}~\ref{0001507} as a definition.
Prove first that $\langle y,x \rangle = \conj{\langle x,y \rangle}$ for all $x$, $y\in V$ and that
$\langle x,x \rangle > 0$ whenever $x \ne 0$.  Next, for arbitrary $z \in V$, define a function
$f\colon V \sto \C\colon x \mapsto \langle x,z \rangle$.  Prove that
   \[ f(x + y) + f(x - y) = 2f(x) \tag{$\ast$} \]
for all $x$, $y \in V$.  Use this to prove that $f(\alpha x) = \alpha f(x)$ for all
$x \in V$ and $\alpha \in \R$.  (Start with $\alpha$ being a natural number.)  Then
show that $f$ is additive. (If $u$ and $v$ are arbitrary elements of $V$ let $x = u + v$
and $y = u - v$.  Use~$(\ast)$.)  Finally prove that $f(\alpha x) = \alpha f(x)$ for
complex $\alpha$ by showing that $f(ix) = i\,f(x)$ for all $x \in V$.  \ns
\end{proof}
















\section{Orthonormal Sets and Bases}\label{sec_onbases}
\begin{defn} A nonempty subset $E$ of a Hilbert space is
 \index{orthonormal}%
\df{orthonormal} if $e \perp f$ for every pair of distinct vectors $e$ and $f$ in $E$ and
$\norm e = 1$ for every $e \in E$.
\end{defn}

\begin{prop} In an inner product space every orthonormal set is linearly independent.
\end{prop}

\begin{prop}[Bessel's inequality]\label{C063141} Let $H$ be a Hilbert space and $E$ be an
 \index{Bessel's inequality}%
orthonormal subset of~$H$.  For every $x \in H$
  \[ \sum_{e \in E} \norm{\langle x,e \rangle}^2 \le \norm x^2\,. \]
\end{prop}

\begin{cor} If $E$ is an orthonormal subset of a Hilbert space $H$ and $x \in H$, then
$\{e \in E\colon \langle x,e \rangle~\ne~0\}$ is countable.
\end{cor}

\begin{prop}\label{C063144} Let $E$ be an orthonormal set in a Hilbert space~$H$. Then the
following are equivalent.
 \begin{enumerate}
  \item $E$ is a maximal orthonormal set.
  \item If $x \perp E$, then $x = 0$ ($E$ is \emph{total}).
 \index{total!orthonormal set}%
  \item $\bigvee E = H$ ($E$ is a \emph{complete orthonormal set}).
 \index{complete!orthonormal set}%
  \item $\D x = \sum_{e \in E} \langle x,e \rangle e$ for all $x \in H$. \emph{(Fourier
expansion)}
 \index{Fourier!expansion}%
  \item $\D \langle x,y \rangle = \sum_{e \in E} \langle x,e \rangle \langle e,y \rangle$
for all $x,y \in H$. \emph{(Parseval's identity.)}
 \index{Parseval's identity}%
  \item $\D \norm{x}^2 = \sum_{e \in E} \abs{\langle x,e \rangle}^2$ for all $x \in H$.
\emph{(Parseval's identity.)}
 \end{enumerate}
\end{prop}

\begin{defn}\label{C063147} If $H$ is a Hilbert space, then an orthonormal set $E$ satisfying
any one (hence all) of the conditions in proposition~\ref{C063144} is called an
 \index{orthonormal!basis}%
 \index{basis!orthonormal}%
\df{orthonormal basis} for~$H$ (or a
 \index{complete!orthonormal set}%
 \index{orthonormal!set!complete}%
\df{complete orthonormal set} in~$H$, or a
 \index{Hilbert!space!basis for}%
 \index{basis!for a Hilbert space}%
\df{Hilbert space basis} for~$H$).  Notice that this is a very different thing from the
usual (Hamel) basis for a vector space (see~\ref{defn_Hamel_basis}).  For a Hilbert space
the basis vectors are required to be of length one and to be mutually perpendicular
(conditions which make no sense in a vector space); but it is \emph{not} necessary that
every vector in the space be a linear combination of the basis vectors.
\end{defn}

\begin{prop} If $E$ is an orthonormal subset of a Hilbert space $H$, then there exists an
orthonormal basis for $H$ which contains~$E$.
\end{prop}

\begin{cor} Every nonzero Hilbert space has a basis.
\end{cor}

\begin{prop}  If $E$ is an orthonormal subset of a Hilbert space $H$ and $x \in H$, then
$\{\langle x,e \rangle e: e \in E\}$ is summable.
\end{prop}

\begin{exam}\label{C063149} For each $n \in \N$ let $\vc e^n$ be the sequence in $l_2$ whose
$n^{\text{th}}$ coordinate is $1$ and all the other coordinates are~$0$. Then $\{\vc e^n\colon n \in \N\}$
is an orthonormal basis for~$l_2$.  This is the
 \index{usual!orthonormal basis for~$l_2$}%
 \index{orthonormal!basis!usual (or standard) for~$l_2$}%
 \index{basis!standard (or usual)}%
\df{usual} (or
 \index{standard!orthonormal basis}%
\df{standard}) \df{orthonormal basis} for~$l_2$.
\end{exam}

\begin{exer} A function $f\colon \R \sto \C$ of the form
   \[ f(t) = a_0 + \sum_{k=1}^n (a_k \cos kt + b_k \sin kt) \]
where $a_0, \dots, a_n, b_1, \dots, b_n$ are complex numbers is a
 \index{trigonometric polynomial}%
 \index{polynomial!trigonometric}%
\df{trigonometric polynomial}.
 \begin{enumerate}
  \item Explain why complex valued $2\pi$-periodic functions on the real line $\R$ are
often identified with complex valued functions on the unit circle~$\T$.
  \item Some authors say that a \emph{trigonometric polynomial} is a function of the form
    \[ f(t) = \sum_{k=-n}^n c_k e^{ikt} \]
where $c_1, \dots, c_n$ are complex numbers. Explain carefully why this is exactly the
same as the preceding definition.
  \item Justify the use of the term \emph{trigonometric polynomial}.
  \item Prove that every continuous $2\pi$-periodic function on $\R$ is the uniform limit
of a sequence of trigonometric polynomials.
 \end{enumerate}
\end{exer}

\begin{prop}\label{prop_unique_onbasis} Let $E$ be a basis for a Hilbert space $H$ and $x \in H$. 
If $x = \sum_{e \in E} \alpha_e e$, then $\alpha_e = \langle x,e \rangle$ for each $e \in E$.
\end{prop}

\begin{prop} Consider the Hilbert space $\lfs 2([0,2\pi],\C)$ of (equivalence classes of complex valued,
 \index{L@$\lfs 2([0,2\pi],\C)$}%
square integrable functions on $[0,2\pi]$ with inner product
  \[ \langle f,g \rangle = \frac1{2\pi}\int_0^{2\pi} f(x) \conj{g(x)}\,dx. \]
For every integer $n$ (positive, negative, or zero) let $\vc e^n(x) = e^{inx}$ for all $x
\in [0,2\pi]$.  Then the set of these functions $\vc e^n$ for $n \in \Z$ is an
orthonormal basis for $\lfs 2([0,2\pi],\C)$.
\end{prop}

\begin{exam} The sum of the infinite series $\D\sum_{k=1}^\infty \frac1{k^2}$ is $\dfrac{\pi^2}6$\,.
\end{exam}

\begin{proof}[\emph{Hint for proof}]   Let $f(x) = x$ for $0 \le x \le 2\pi$.  Regard $f$ as a member of
$\lfs 2([0,2\pi])$.  Find $\norm{f}$ in two ways.  \ns
\end{proof}

The next proposition allows us to define the notion of the \emph{dimension} of a Hilbert space.  (In finite dimensional
spaces it agrees with the vector space concept.)

\begin{prop} Any two orthonormal bases of a Hilbert space have the same cardinality (that is, they are in
one-to-one correspondence).
\end{prop}

\begin{proof}  See \cite{Halmos:1951}, page 29, Theorem 1.   \ns
\end{proof}

\begin{defn} The
 \index{dimension}%
\df{dimension} of a Hilbert space is the cardinality of any orthonormal basis for the space.
\end{defn}

\begin{prop}\label{prop_Schauder_basis_Hsp} A Hilbert space is separable if and only if it has a
 \index{separable!Hilbert space}%
 \index{Hilbert!space!separable}%
countable orthonormal basis.
\end{prop}

\begin{proof}  See \cite{Kreyszig:1978}, page 171, Theorem 3.6-4.   \ns
\end{proof}





















\section{The Riesz-Fr\'echet Theorem}
\begin{exam}\label{000341} Let $H$ be a Hilbert space and $a \in H$.  Define $\psi_a\colon
H \sto \C\colon x \mapsto \langle x,a \rangle$.  Then $\psi_a \in H^*$.
\end{exam}

Now we generalize theorem~\ref{0001601} to the infinite dimensional setting. This result
is sometimes called the \emph{Riesz representation theorem} (which invites confusion with
 \index{representation!of continuous linear functionals}%
 \index{representation!Riesz's little theorem}%
the more substantial result about representing  certain linear functionals as measures)
or the \emph{little Riesz representation theorem}.  This theorem says that the \emph{only}
continuous linear functionals on a Hilbert space are the functions $\psi_a$ defined in~\ref{000341}.
In other words, given a continuous linear functional $f$ on a Hilbert space there is a unique vector
$a$ in the space such that the action of $f$ on a vector $x$ can be represented simply by taking the
inner product of $x$ with~$a$.

\begin{thm}[Riesz-Fr\'echet Theorem]\label{000342} If $f \in H^*$ where $H$ is a Hilbert
 \index{Riesz-Fr\'echet theorem}%
space, then there exists a unique vector $a$ in $H$ such that $f = \psi_a$. Furthermore,
$\norm a = \norm{\psi_a}$.
\end{thm}

\begin{proof}[\emph{Hint for proof}] For the case $f \ne 0$, choose a unit vector $z$ in $(\ker f)^\perp$.
Notice that for every $x \in H$ the vector $x - \frac{f(x)}{f(z)}\,z$ belongs to $\ker f$.  \ns
\end{proof}

\begin{cor} The kernel of every nonzero bounded linear functional on a Hilbert space has codimension~$1$.
\end{cor}

\begin{proof}[\emph{Hint for proof}] The
 \index{codimension}%
\df{codimension} of a subspace of a Hilbert space is the dimension of its orthogonal complement.  This is really
a corollary to the \emph{proof} of the \emph{Riesz-Fr\'echet theorem}.  In our proof we found that the vector
representing the bounded linear functional $f$ was a (nonzero) multiple of a unit vector in $(\ker f)^\perp$.
What would happen if there were two linearly independent vectors in $(\ker f)^\perp$?   \ns
\end{proof}

\begin{defn} Recall that a mapping $T\colon V \sto W$ between complex vector spaces is
 \index{conjugate linear}%
 \index{linear!conjugate}%
\df{conjugate linear} if $T(x + y) = Tx + Ty$ and $T(\alpha x) = \conj{\alpha}\,Tx$ for all
$x$, $y \in V$ and $\alpha \in \C$. A bijective conjugate linear map is an
 \index{anti-isomorphism}%
 \index{isomorphism!anti-}%
\df{anti-isomorphism}.
\end{defn}

\begin{exam}\label{0022057} For a Hilbert space $H$, the mapping
$\psi\colon H \sto H^*\colon a \mapsto \psi_a$ defined in example~\ref{000341} is an
isometric anti-isomorphism between $H$ and its dual space.
\end{exam}

\begin{defn} A conjugate linear mapping $C\colon V \sto V$ from a vector space into itself
which satisfies $C^2 = \id V$ is called a
 \index{conjugation}
\df{conjugation} on~$V$.
\end{defn}

\begin{exam} Complex conjugation $z \mapsto \conj z$ is an example of a conjugation in the
vector space~$\C$.
\end{exam}

\begin{exam}\label{0022067} Let $H$ be a Hilbert space and $E$ be a basis for~$H$.  Then the
map $x \mapsto \sum_{e \in E}\langle e,x \rangle e$ is a conjugation and an isometry on~$H$.
\end{exam}

The term ``anti-isomorphism'' is a bit misleading. It suggests something entirely
different from an isomorphism.  In fact, an anti-isomorphism is nearly as good as an
isomorphism. The next proposition says in essence that a Hilbert space and its dual are
isomorphic \emph{because} they are anti-isomorphic.

\begin{prop}\label{0022071} Let $H$ be a Hilbert space, $\psi\colon H \sto H^*$ be the
anti-isomorphism defined in~\ref{000341} (see~\ref{0022057}), and $C$ be the conjugation
defined in~\ref{0022067}. Then the composite $C \psi^{-1}$ is an isometric isomorphism
from $H^*$ onto~$H$.
\end{prop}

\begin{cor}\label{cor_Hsp_iso_dual} Every Hilbert space is isometrically isomorphic to its dual.
\end{cor}

\begin{prop} Let $H$ be a Hilbert space and $\psi\colon H \sto H^*$ be the anti-isomorphism
defined in~\ref{000341} (see~\ref{0022057}).  If we define $\langle f,g \rangle :=
\langle \psi^{-1}g, \psi^{-1}f \rangle$ for all $f$, $g \in H^*$, then $H^*$ becomes a
Hilbert space isometrically isomorphic to~$H$.  The resulting norm on $H^*$ is its usual
norm.
\end{prop}

\begin{cor} Every Hilbert space is reflexive.
\end{cor}

\begin{exam} Let $H$ be the collection of all absolutely continuous complex
valued functions $f$ on $[0,1]$ such that $f(0) = 0$ and $f' \in
\lfs 2([0,1],\lambda)$.  Define an inner product on $H$ by
 \[\langle f,g \rangle := \int_0^1 f'(t)\,\clo{g'(t)}\,dt.\]
We already know that $H$ is a Hilbert space (see example~\ref{exam_Hsp_abs_cont})
Fix $t \in (0,1]$.  Define
   \[ E_t\colon H \sto \C\colon f \mapsto f(t). \]
Then $E_t$ is a bounded linear functional on~$H$. What is its $\norm{E_t}$?
What vector $g$ in $H$ represents the functional~$E_t$?
\end{exam}

\begin{exer} One afternoon your pal Fred R. Dimm comes to you with a problem. ``Look,'' he says,
``On $\lfs 2 = \lfs 2(\R,\R)$ the evaluation functional $E_0$, which evaluates each member of $\lfs 2$
at~$0$, is clearly a bounded linear functional.  So by the Riesz representation theorem there
should be some function $g$ in $\lfs 2$ such that $f(0) = \langle f,g \rangle = \int_{-\infty}^\infty fg$
for all $f$ in~$\lfs 2$.  But that's a property of the `$\delta$-function', which I was told
doesn't exist.  What am I doing wrong?'' Give Freddy some (helpful) advice.
\end{exer}

\begin{exer} Later the same afternoon Freddy's back.  After considering the
advice you gave him (in the preceding problem) he has revised his
question by letting the domain of $E_0$ be the set of bounded
continuous real valued functions defined on~$\R$.  Patiently you
explain to Freddy that there are now at least two things wrong his
invocation of the `$\delta$-function'.  What are they?
\end{exer}

\begin{exer} The nightmare continues. It's 11 P. M. and Freddy's back again.
This time (happily having given up on $\delta$-functions) he wants
to apply the representation theorem to the functional ``evaluation
at zero'' on the space~$l_2(\Z)$.  Can he do it?  Explain.
\end{exer}



























\section{Strong and Weak Topologies on Hilbert Spaces}
We recall briefly the definition and crucial facts about \emph{weak topologies}.  For a more detailed introduction
to this class of topologies see any good textbook on topology or section 11.4 of my notes~\cite{Erdman:2007}.

\begin{defn}\label{C021414} Suppose that $S$ is a set, that for every $\alpha \in A$ (where
$A$ is an arbitrary index set) $X_\alpha$ is a topological space, and that $f_\alpha
\colon S \sto X_\alpha$ for every $\alpha \in A$. Let
   \[ \sfml S = \{{f_\alpha}^\gets(U_\alpha) \colon \open{U_\alpha}{X_\alpha}
                                             \text{ and } \alpha \in A\}\,. \]
Use the family $\sfml S$ as a subbase for a topology on~$S$.  This topology is called the
 \index{weak!topology!induced by a family of functions}%
 \index{topology!weak!induced by a family of functions}%
\df{weak topology} induced by (or determined by) the functions~$f_\alpha$.
\end{defn}

\begin{prop}\label{C021421} Under the weak topology (defined above) on a set $S$ each of the functions $f_\alpha$ is continuous;
in fact the weak topology is  the weakest topology on
 \index{weak!topology!characterization of}%
 \index{topology!weak!characterization of}%
$S$ under which these functions are continuous.
\end{prop}

\begin{prop}\label{C021424}  Let $X$ be a topological space with the weak topology determined
by a family $\fml F$ of functions. Prove that a function $g\colon W \sto X$, where $W$ is
a topological space, is continuous if and only if $f \circ g$ is continuous for every $f
\in \fml F$.
\end{prop}

\begin{defn}\label{C021437}  Let $\bigl(A_\lambda\bigr)_{\lambda \in \Lambda}$ be an indexed
family of topological spaces. The weak topology on the Cartesian product $\prod
A_\lambda$ induced by the family of coordinate projections $\pi_\lambda$ (see
definition~\ref{C015531}) is called the
 \index{product!topology}%
 \index{topology!product}%
\df{product topology}.
\end{defn}

\begin{prop} The product of a family of Hausdorff spaces is Hausdorff.
\end{prop}

\begin{exam}\label{C031473} Let $Y$ be a topological space with the weak topology (see
definition~\ref{C021414}) determined by an indexed family $(f_\alpha)_{\alpha \in A}$ of
functions where for each $\alpha \in A$ the codomain of $f_\alpha$ is a topological
space~$X_\alpha$. Then a net $\bigl(y_{{}_\lambda}\bigr)_{\lambda \in \Lambda}$ in $Y$
converges to a point $a \in Y$ if and only if
\mbox{$f_\alpha(y_{{}_\lambda})~\to_{\lambda \in \Lambda}~f_\alpha(a)$} for every $\alpha
\in A$.
\end{exam}

\begin{exam}\label{C031475} If $Y = \prod_{\alpha \in A} X_\alpha$ is a product of nonempty
topological spaces with the product topology (see definition~\ref{C021437}), then a net
$\bigl(\sbsb y\lambda\bigr)$ in $Y$ converges to $a \in Y$ if and only if $\bigl(\sbsb
y\lambda\bigr)_\alpha \sto a_\alpha$ for every $\alpha \in A$.
\end{exam}

\begin{exam}\label{exam_prod_top_norm} If $V$ and $W$ are normed linear spaces, then the product topology on
$V \times W$ is the same as the topology induced by the product norm on $V \oplus W$.
\end{exam}

The next example illustrates the inadequacy of sequences when dealing with general topological spaces.

\begin{exam} Let $\fml F = \fml F([0,1])$ be the set of all functions $f\colon [0,1] \sto \R$ and give
$\fml F$ the product topology, that is, the weak topology determined by the evaluation
functionals
    \[ E_x\colon \fml F \sto [0,1]\colon f \mapsto f(x) \]
where $x \in [0,1]$.  Thus a basic open neighborhood of a point $g \in \fml F$ is determined
by a finite set of points $A \in \fin[0.1]$ and a number $\epsilon > 0$:
  \[ U(g\,;\,A\,;\, \epsilon) := \{f \in \fml F\colon \abs{f(t) - g(t)} < \epsilon
                           \text{ for all $t \in A$}\}. \]
(What is usual name for this topology?) Let $\fml G$ be the set of all functions in $\fml F$
having finite support. Then
  \begin{enumerate}
    \item the constant function $\vc 1$ belongs to $\clo{\fml G}$,
    \item there is a net of functions in $\fml G$ which converges to~$\vc 1$, but
    \item no sequence in $\fml G$ converges to~$\vc 1$.
  \end{enumerate}
\end{exam}

\begin{defn} A net $(x_{{}_{\sst\lambda}})$ in a Hilbert space $H$ is said to
 \index{convergence!in Hilbert space!weak}%
 \index{weak!convergence!in Hilbert space}%
\df{converge weakly} to a point $a$ in $H$ if $\langle x{{}_{\sst\lambda}}, y \rangle \to \langle a,y \rangle$
for every $y \in H$. In this case we write $x_{{}_{\sst\lambda}} \to^w a$. In a Hilbert space a net is said to
 \index{convergence!in Hilbert space!strong}%
 \index{strong!convergence!in Hilbert space}%
\df{converge strongly} (or
 \index{convergence!in norm}%
 \index{norm!convergence}%
\df{converge in norm}) if it converges with respect to the topology induced by the norm. This is the usual type of
convergence in a Hilbert space. If we wish to emphasize a distinction between modes of convergence we may write
$x_{{}_{\sst\lambda}} \to^s a$ for strong convergence.
\end{defn}

\begin{exer} Explain why the topology on a Hilbert space generated by the weak convergence of nets is in fact a
weak topology in the usual topological sense of the term.
\end{exer}

When we say that a subset of a Hilbert space $H$ is
 \index{weak!closure}%
 \index{closed!weakly}%
\df{weakly closed} we mean that it is closed with respect to the weak topology on $H$; a set is
 \index{weak!compactness}%
 \index{compact!weakly}%
\df{weakly compact} if it is compact in the weak topology on~$H$; and so on. A function
$f\colon H \sto H$ is
 \index{weak!continuity}%
 \index{continuous!weakly}%
\df{weakly continuous} if it is continuous as a map from $H$ with the weak topology into $H$
with the weak topology.

\begin{exer} Let $H$ be a Hilbert space and $(x_{{}_{\sst\lambda}})$ be a net in $H$.
 \begin{enumerate}
  \item If $x_{{}_{\sst\lambda}} \to^s  a$ in $H$, then $x_{{}_{\sst\lambda}} \to^w a$ in $H$
  \item The strong (norm) topology on $H$ is stronger (bigger) than the weak topology.
  \item Show by example that the converse of (a) does not hold.
  \item Is the norm on $H$ weakly continuous?
  \item If $x_{{}_{\sst\lambda}} \to^w  a$ in $H$ and if $\norm{x_{{}_{\sst\lambda}}} \to \norm{a}$, then
$x_{{}_{\sst\lambda}}\to^s a$ in~$H$.
  \item If a linear map $T\colon H \sto H$ is (strongly) continuous, then it is weakly continuous.
 \end{enumerate}
\end{exer}

\begin{exer} Let $H$ be an infinite dimensional Hilbert space. Consider the following subsets of~$H$:
 \begin{enumerate}[label=\rm{(\alph*)}]
  \item the open unit ball;
  \item the closed unit ball;
  \item the unit sphere;
  \item a closed linear subspace.
 \end{enumerate}
Determine for each of these sets whether it is each of the following: strongly closed, weakly
closed, strongly open, weakly open, strongly compact, weakly compact.  (One part of this is
too hard at the moment: \emph{Alaoglu's theorem}~\ref{C067441} will tell us that the
closed unit ball is weakly compact.)
\end{exer}

\begin{exer} Let $\{e^n\colon n \in \N\}$ be the usual basis for $l_2$. For
each $n \in \N$ let $a_n = \sqrt n\, e^n$.  Which of the following
are correct?
 \begin{enumerate}
  \item $0$ is a weak accumulation point of the set
$\{a_n\colon n \in \N\}$.
  \item $0$ is a weak cluster point of the sequence~$(a_n)$.
  \item $0$ is the weak limit of a subsequence of~$(a_n)$.
 \end{enumerate}
\end{exer}










\section{Universal Morphisms}
Much of mathematics involves the construction of new objects from old ones---things such
as products, coproducts, quotients, completions, compactifications, and unitizations.
Often it is possible---and highly desirable---to characterize such a construction by
means of a diagram which describes what the constructed object ``does'' rather than
telling what it ``is'' or how it is constructed.  Such a diagram is a
 \index{universal!mapping!diagram}%
 \index{mapping!universal}%
\df{universal mapping diagram} and it describes the
 \index{universal!mapping!property}%
 \index{mapping!diagram!universal}%
\df{universal property} of the object being constructed.
In particular it is usually possible to characterize such a construction by the existence
of a unique morphism having some particular property. Because this morphism and its
corresponding property uniquely characterize the construction in question, they are referred to as a
\emph{universal morphism} and a \emph{universal property}, respectively. The following definition
is one way of describing the action of such a morphism.  If this is your first meeting with the
concept of \emph{universal} don't be alarmed by its rather abstract nature.  Few people, I think,
feel comfortable with this definition until they have encountered a half dozen or more examples in
different contexts.  (There exist even more technical definitions for this important concept. For
example, \emph{an initial or terminal object in a comma category}. To unravel this, if you are
interested, search for ``universal property'' in Wikipedia~\cite{wiki:xxx}.)

\begin{defn}\label{00078} Let $\cat A \to^{\ftr{F}} \cat B$ be a functor between categories
$\cat A$ and $\cat B$ and $B$ be an object in~$\cat B$.  A pair $(A,u)$ with $A$ an
object in $\cat A$ and $u$ a $\cat B$-morphism from $B$ to~$\ftr F(A)$ is a
 \index{universal!morphism}%
 \index{morphism!universal}%
\df{universal morphism} for $B$ (with respect to the functor~$\ftr F$) if for every
object $A'$ in $\cat A$ and every $\cat B$-morphism $B \to^f \ftr F(A')$ there exists a
unique $\cat A$-morphism $A \to^{\tilde f} A'$ such that the following diagram commutes.
 \begin{equation}\label{00078i}
   \xy
     \qtriangle/{>}`{>}`{-->}/[B`\ftr F(A)`\ftr F(A');u`f`\ftr F(\tilde f)]
     \morphism(1000,500)|r|/{-->}/<0,-500>[A`A';\tilde f]
   \endxy
 \end{equation}
In this context the object $A$ is often referred to as a
 \index{universal!object}%
 \index{object!universal}%
\df{universal object} in~$\cat A$.
\end{defn}

Here is a first example of such a property.

\begin{defn} Let $F$ be an object in a concrete category $\cat C$ and
$\iota\colon S \sto \abs F$ be an inclusion map whose domain is a nonempty set~$S$.  We say that the
object $F$ is
 \index{free!object}%
 \index{object!free}%
\df{free on} the set $S$ (or that $F$ is the \df{free object generated by}~$S$) if for
every object $A$ in $\cat C$ and every map $f\colon S \sto \abs A$ there exists a unique
morphism $\widetilde f_\iota\colon F \sto A$ in $\cat C$ such that $\abs{\wt f_\iota}
\circ \iota = f$.
 \[\xy
     \qtriangle/{>}`{>}`{-->}/[S`\abs F`\abs A;\iota`f`\abs{\wt f_\iota}]
     \morphism(1000,500)|r|/{-->}/<0,-500>[F`A;\widetilde f_\iota]
   \endxy\]
We will be interested in
 \index{free!vector space}%
 \index{vector!space!free}%
\df{free vector spaces}; that is, free objects in the category $\cat{VEC}$ of vector
spaces and linear maps.  Naturally, merely \emph{defining} a concept does not guarantee
its existence. It turns out, in fact, that free vector spaces exist on arbitrary sets.
(See exercise~\ref{ump001z}.)
\end{defn}

\vskip 5 pt

\begin{exer} In the preceding definition reference to the forgetful functor is often
omitted and the accompanying diagram is often drawn as follows:
 \[\xy
     \qtriangle/{>}`{>}`{-->}/[S`F`A;\iota`f`\wt f_\iota]
   \endxy\]
It certainly looks a lot simpler.  Why do you suppose I opted for the more complicated
version?
\end{exer}

\vskip 5 pt

\begin{prop} If two objects in some concrete category are free on the same set,
then they are isomorphic.
\end{prop}

\vskip 5 pt

\begin{defn} Let $A$ be a subset of a nonempty set~$S$ and $\K$ be either $\R$ or~$\C$. Define
$\sbsb\chi A \colon S \sto \K$, the
 \index{characteristic!function}%
 \index{<@$\sbsb{\chi}{A}$ (characteristic function of $A$)}%
 \index{function!characteristic}%
\df{characteristic function} of~$A$, by
  \[ \sbsb{\chi}{A}(x) = \begin{cases}
                                  1, &\text{if $x \in A$} \\
                                  0, &\text{otherwise}
                              \end{cases} \]
\end{defn}

\begin{exam}\label{ump001z} If $S$ is an arbitrary nonempty set,
then there exists a vector space $V$ over~$\K$ which is free on~$S$.  This vector space
is unique (up to isomorphism).
\end{exam}

\begin{proof}[\emph{Hint for proof}] Given the set $S$ let $V$ be the set of all
$\K$-valued functions on $S$ which have finite support. Define addition and scalar
multiplication pointwise. The map $\iota\colon s \mapsto \sbsb \chi {\{s\}}$ of each
element $s \in S$ to the characteristic function of~$\{s\}$ is the desired injection. To
verify that $V$ is free over $S$ it must be shown that for every vector space $W$ and
every function \hbox{$S~\to^{f}~\abs W$} there exists a unique linear map $V \to^{\wt f}
W$ which makes the following diagram commute.
 \[\xy
 \index{l@$l_2 = l_2(\N)$!square summable sequences}%
 \index{l@$l_2 = l_2(\N)$!as an inner product space}%
 \index{inner product!space!$l_2$ as a}%
     \qtriangle/{>}`{>}`{-->}/[S`\abs V`\abs W;\iota`f`\abs{\wt f}]
     \morphism(1100,500)|r|/{-->}/<0,-500>[V`W;\wt f]
   \endxy\]
\ns \end{proof}

\begin{prop} Every vector space is free.
\end{prop}

\begin{proof}[\emph{Hint for proof}] Of course, part of the problem is to specify a set
$S$ on which the given vector space is free. \ns
\end{proof}


\begin{exam} Let $S$ be a nonempty set and let $S' = S \cup \{\vc 1\}$ where $\vc 1$ is any
element not belonging to~$S$.  A
 \index{word}%
\df{word} in the language~$S$ is a sequence $s$ of elements of $S'$ which is eventually
$\vc 1$ and satisfies: if $s_k = \vc 1$, then $s_{k+1} = \vc 1$.  The constant sequence
$(\vc 1, \vc 1, \vc 1, \dots)$ is called the
 \index{empty word}%
 \index{word!empty}%
\df{empty word}. Let $F$ be the set of all words of members in the language~$S$. Suppose
that $s = (s_1, \dots, s_m, \vc 1, \vc 1, \dots)$ and $t = (t_1, \dots, t_n, \vc 1, \vc
1, \dots)$ are words (where $s_1$, \dots, $s_m$, $t_1$, \dots, $t_n \in S$).  Define
   \[ x \ast y := (s_1, \dots, s_m, t_1, \dots, t_n, \vc 1, \vc 1, \dots). \]
This operation is called
 \index{concatenation}%
\df{concatenation}. It is not difficult to see that the set $F$ under concatenation is a
 \index{monoid}%
monoid (a semigroup with identity) where the empty word is the identity element.  This is the
 \index{free!monoid}%
 \index{monoid!free}%
\df{free monoid} generated by~$S$.  If we exclude the empty word we have the
 \index{free!semigroup}%
 \index{semigroup!free}%
\df{free semigroup} generated by~$S$. The associated diagram is
 \[\xy
     \qtriangle/{>}`{>}`{-->}/[S`\abs F`\abs G;\iota`f`\abs{\tilde f}]
     \morphism(1100,500)|r|/{-->}/<0,-500>[F`G;\tilde f]
   \endxy\]
where $\iota$ is the obvious injection $s \mapsto (s,\vc1, \vc 1, \dots)$ (usually
treated as an inclusion mapping), $G$ is an arbitrary semigroup, $f\colon S \sto G$ is an
arbitrary function, and $\abs{\hphantom{O}}$ is the forgetful functor from the category of monoids and
homomorphisms (or the category of semigroups and homomorphisms) to the category
$\cat{SET}$.
\end{exam}

\begin{exam} Recall from definition~\ref{C015611} the usual presentation of the coproduct of two objects in a category. If
$A_1$ and $A_2$ are two objects in a category $\cat C$, then a
 \index{coproduct}%
\df{coproduct} of $A_1$ and $A_2$ is a triple $(Q,\iota_1,\iota_2)$ with $Q$ an object in
$\cat C$ and $A_k \to^{\iota_k} Q$ ($k = 1,2$) morphisms in $\cat C$ which satisfies the
following condition: if $B$ is an arbitrary object in $\cat C$ and $A_k \to^{f_k} B$ ($k
= 1$,$2$) are arbitrary morphisms in~$\cat C$, then there exists a unique $\cat C$-morphism
$Q \to^f B$ which makes the following diagram commute.
  \begin{equation}
    \xy
      \Atrianglepair/<-`<-`<-`>`<-/[B`A_1`Q`A_2;f_1`f`f_2`\iota_1`\iota_2]
    \endxy
   \end{equation}

It may not be obvious at first glance that this construction is \emph{universal} in the
sense of definition~\ref{00078}.  To see that it in fact is, let $\ftr D$ be the diagonal
functor from a category $\cat C$ to the category of pairs $\cat{C^2}$ (see
example~\ref{0007606}). Suppose that $(Q,\iota_1,\iota_2)$ is a coproduct of the $\cat
C$-objects $A_1$ and $A_2$ in the sense defined above.  Then $A = (A_1,A_2)$ is an object
in $\cat{C^2}$, $A \to^{\iota} \ftr D(Q)$ is a $\cat{C^2}$-morphism, and the pair
$(A,\iota)$ is universal in the sense of~\ref{00078}.  The diagram corresponding to
diagram~\eqref{00078i} is
 \begin{equation}
   \xy
     \qtriangle/{>}`{>}`{-->}/[A`\ftr D(Q)`\ftr D(B);\iota`f`\ftr D(\tilde f)]
     \morphism(1000,500)|r|/{-->}/<0,-500>[Q`B;\tilde f]
   \endxy
 \end{equation}
where $B$ is an arbitrary object in $\cat C$ and (for $k = 1,2$) $A_k \to^{f_k} B$ are
arbitrary $\cat C$-morphisms so that $f = (f_1,f_2)$ is a $\cat{C^2}$-morphism.
\end{exam}

\begin{exam} The coproduct of two objects $H$ and $K$ in the category $\cat{HIL}$ of
 \index{coproduct!in $\cat{HIL}$}%
 \index{hil@$\cat{HIL}$!coproducts in}%
Hilbert spaces and bounded linear maps (and more generally in the category of inner
product spaces and linear maps) is their (external orthogonal) direct sum $H \oplus K$
(see~\ref{0001504}).
\end{exam}


\begin{exam}\label{00078036} Let $A$ and $B$ be Banach spaces. On the Cartesian
 \index{coproduct!in $\cat{BAN_\infty}$}%
 \index{ban@$\cat{BAN_\infty}$!coproducts in}%
product $A \times B$ define addition and scalar multiplication pointwise. For every
$(a,b) \in A \times B$ let $\norm{(a,b)} = \max\{\norm a,\norm b\}$. This makes $A \times
B$ into a Banach space, which is denoted by $A \oplus B$ and is called the
 \index{direct!sum!of Banach spaces}%
\df{direct sum} of $A$ and~$B$.  The direct sum is a coproduct in the topological
category $\cat{BAN_\infty}$ of Banach spaces but not in the corresponding geometrical
category~$\cat{BAN_1}$.
\end{exam}

\begin{exam}\label{00078037} To construct a coproduct on the geometrical category
 \index{coproduct!in $\cat{BAN_1}$}%
 \index{ban@$\cat{BAN_1}$!coproducts in}%
$\cat{BAN_1}$ of Banach spaces define the vector space operations pointwise on $A \times
B$ but as a norm use $\norm{(a,b)}_1 = \norm a + \norm b$ for all $(a,b) \in A \times B$.
\end{exam}

Virtually everything in category theory has a dual concept---one that is obtained by
reversing all the arrows. We can, for example, reverse all the arrows in
diagram~\eqref{00078i}.


\begin{defn}\label{000781} Let $\cat A \to^{\ftr{F}} \cat B$ be a functor between categories
$\cat A$ and $\cat B$ and $B$ be an object in~$\cat B$.  A pair $(A,u)$ with $A$ an
object in $\cat A$ and $u$ a $\cat B$-morphism from $\ftr F(A)$ to~$B$  is a
 \index{co-universal!morphism}%
 \index{morphism!co-universal}%
\df{co-universal morphism for $B$ (with respect to}~$\ftr F$) if for every object $A'$ in
$\cat A$ and every $\cat B$-morphism $\ftr F(A') \to^f B$ there exists a unique $\cat
A$-morphism $A' \to^{\tilde f} A$ such that the following diagram commutes.
 \begin{equation}\label{000781i}
   \xy
     \qtriangle/{<-}`{<-}`{<--}/[B`\ftr F(A)`\ftr F(A');u`f`\ftr F(\tilde f)]
     \morphism(1000,500)|r|/{<--}/<0,-500>[A`A';\tilde f]
   \endxy
 \end{equation}
Some authors reverse the convention and call the morphism in~\ref{00078}
\emph{co-universal} and the one here \emph{universal}.  Other authors, this one included,
call both \emph{universal morphisms}.
\end{defn}

\begin{conv} Morphisms of both the types defined in~\ref{00078} and~\ref{000781} will
 \index{conventions!about universality}%
be referred to as \emph{universal morphisms}.
\end{conv}

\begin{exam}\label{0007812} Recall from definition~\ref{C015511} the categorical approach to products.  If
$A_1$ and $A_2$ are two objects in a category $\cat C$, then a
 \index{product}%
\df{product} of $A_1$ and $A_2$ is a triple $(P,\pi_1,\pi_2)$ with $P$ an object in $\cat
C$ and $Q \to^{\iota_k} A_k$ ($k = 1,2$) morphisms in $\cat C$ which satisfies the
following condition: if $B$ is an arbitrary object in $\cat C$ and $B \to^{f_k} A_k$ ($k
= 1,2$) are arbitrary morphisms in~$\cat C$, then there exists a unique $\cat C$-morphism
$B \to^f P$ which makes the following diagram commute.
  \begin{equation}
    \xy
      \Atrianglepair/>`>`>`<-`>/[B`A_1`P`A_2;f_1`f`f_2`\pi_1`\pi_2]
    \endxy
   \end{equation}
This is (co)-universal in the sense of definition~\ref{000781}.
\end{exam}

\begin{exam} The product of two objects $H$ and $K$ in the category $\cat{HIL}$ of
 \index{product!in $\cat{HIL}$}%
 \index{hil@$\cat{HIL}$!products in}%
Hilbert spaces and bounded linear maps (and more generally in the category of inner
product spaces and linear maps) is their (external orthogonal) direct sum $H \oplus K$
(see~\ref{0001504}).
\end{exam}


\begin{exam} If $A$ and $B$ are Banach algebras their direct sum $A \oplus B$ is a
 \index{product!in $\cat{BA_\infty}$}%
 \index{ba@$\cat{BA_\infty}$!products in}%
 \index{product!in $\cat{BA_1}$}%
 \index{ba@$\cat{BA_1}$!products in}%
product in both the topological and geometric categories, $\cat{BA_\infty}$ and
$\cat{BA_1}$, of Banach algebras.  Compare this to the situation discussed in
examples~\ref{00078036} and~\ref{00078037}.
\end{exam}

Universal constructions in any category whatever produce objects that are unique up to isomorphism.

\begin{prop}\label{C035562} Universal objects in a category are
 \index{universal!object!uniqueness of}%
essentially unique.
\end{prop}















\section{Completion of an Inner Product Space}
A little review of the topic of completion of metric spaces may be useful here.  Occasionally
it may be tempting to think that the completion of a metric space $M$ is a complete metric space
of which some dense subset is homeomorphic to $M$.  The problem with this characterization is that
it may very well produce two ``completions'' of a metric space which are not even homeomorphic.

\begin{exam} Let $M$ be the interval $(0,\infty)$ with its usual metric. There exist two complete
metric spaces $N_1$ and $N_2$ such that
  \begin{enumerate}
    \item $N_1$ and $N_2$ are \emph{not} homeomorphic,
    \item $M$ is homeomorphic to a dense subset of $N_1$, and
    \item $M$ is homeomorphic to a dense subset of~$N_2$.
  \end{enumerate}
\end{exam}

Here is the correct definition.

\begin{defn}\label{C035511} Let $M$ and $N$ be metric spaces.  We say that $N$ is a
 \index{completion!of a metric space}%
 \index{metric!space!completion of a}%
\df{completion} of $M$ if $N$ is complete and $M$ is isometric to a dense subset of~$N$.
\end{defn}

In the construction of the real numbers from the rationals one approach is to treat the
set of rationals as a metric space and define the set of reals to be its completion. One
standard way of producing the completion involves equivalence classes of Cauchy sequences
of rational numbers. This technique extends to metric spaces: an elementary way of
completing an arbitrary metric space starts with defining an equivalence relation on the
set of Cauchy sequences of elements of the space.  Here is a slightly less elementary but
considerably simpler approach.

\begin{prop} Every metric space $M$ is isometric to subspace of~$\fml B(M)$.
\end{prop}

\begin{proof}[\emph{Hint for proof}] If $(M,d)$ is a metric space fix $a \in M$. For each
$x \in M$ define $\phi_x\colon M \sto \R$ by $\phi_x(u) = d(x,u) - d(u,a)$.  Show first
that $\phi_x \in \fml B(M)$ for every $x \in M$.  Then show that $\phi\colon M \sto \fml
B(M)$ is an isometry. \ns
\end{proof}

\begin{cor}\label{C035521} Every metric space has a completion.
\end{cor}

The next proposition shows that the completion of a metric space is universal in the sense of
definition~\ref{00078}.  In this result the categories of interest are the category of all metric spaces
and uniformly continuous maps and its subcategory consisting of all complete metric spaces and uniformly
continuous maps.  Here the forgetful functor $\abs{\hphantom{O}}$ forgets only about completeness of objects
(and does not alter morphisms).

\begin{prop}\label{C035557} Let $M$ be a metric space and $\clo M$ be its completion. If $N$
is a complete metric space and $f\colon M \sto N$ is uniformly continuous, then there
exists a unique uniformly continuous function $\clo f\colon \clo M \sto N$ which makes
the following diagram commute.
 \[\xy
     \qtriangle/{>}`{>}`{-->}/[M`\abs{\clo M}`\abs N;\iota`f`\abs{\clo f}]
     \morphism(1100,500)|r|/{-->}/<0,-500>[\clo M`N;\clo f]
   \endxy\]
\end{prop}

The following consequence of propositions~\ref{C035557} and~\ref{C035562} allows us to
speak of \emph{the} completion of a metric space.

\begin{cor}\label{C035564} Metric space completions are unique (up to isometry).
\end{cor}

Now, what about the completion of inner product spaces?  As we have seen, every inner product space is a
metric space; so it has a (metric space) completion.  The question is: Is this completion a Hilbert space?
The answer is, happily, \emph{yes}.

\begin{thm} Let $V$ be an inner product space and $H$ be its metric space completion.  Then there exists
an inner product on $H$ which
  \begin{enumerate}
    \item is an extension of the one on~$V$ and
    \item induces the metric on~$H$.
  \end{enumerate}
\end{thm}




















\endinput
